% Capitulo: Conhecer o negócio
\chapter{Conhecer o negócio}
\label{cap:negocio}

\rascunho{
  \textbf{Ideia do capítulo.} Com a IA acelerando a parte técnica, o que separa um bom
  profissional de dados de um excelente é o conhecimento de negócio: a capacidade de
  antecipar necessidades e entregar o que ainda não foi pedido, mas que o negócio
  precisava. Isso vem de experiência e de estudo numa área específica.
  \begin{itemize}
    \item O argumento: quando escrever a consulta, o modelo ou o dashboard ficou mais
      rápido, o gargalo passa a ser saber \emph{o que} construir. A IA responde bem
      o que foi perguntado; ela não sabe o que deixou de ser perguntado.
    \item Bom contra excelente: o bom entrega o pedido com qualidade; o excelente
      percebe a pergunta por trás do pedido e chega com a resposta para a próxima.
    \item Profundidade num domínio: entender como a empresa ganha dinheiro, quais
      números o gestor olha toda segunda-feira, onde a operação costuma quebrar.
    \item Relatos da minha experiência em três áreas: agricultura, finanças e
      logística. Um exemplo de cada de necessidade antecipada (algo entregue antes de
      ser pedido). \pendente{escolher os três casos}
    \item Como estudar um domínio de negócio: conversar com a operação, ler os
      relatórios que a diretoria lê, aprender o vocabulário, acompanhar o ciclo do
      negócio (safra, fechamento contábil, sazonalidade de frete).
    \item Ligação com o capítulo~\ref{cap:testes}: quem conhece o negócio escreve os
      testes certos antes do modelo, porque já sabe quais regras não podem quebrar.
    \item Ligação com a Parte~\ref{parte:agil}: antecipar necessidade também é
      refinamento; o profissional que conhece o negócio traz a história pronta para
      o quadro.
  \end{itemize}
}

\section{O gargalo mudou de lugar}
\label{sec:gargalo}

\pendente{escrever}

\section{Do bom ao excelente: entregar o que não foi pedido}
\label{sec:bom-excelente}

\pendente{escrever}

\section{Três domínios: agricultura, finanças e logística}
\label{sec:tres-dominios}

\pendente{escrever, com um relato por domínio}

\section{Como estudar um domínio de negócio}
\label{sec:estudar-dominio}

\pendente{escrever}
